\documentclass[11pt]{article}

\usepackage{amsmath,amssymb,amsthm,bm,mathtools}
\usepackage{geometry}
\geometry{margin=1in}
\usepackage{setspace}
\usepackage{enumitem}
\usepackage{booktabs}
\usepackage{hyperref}
\usepackage{algorithm}
\usepackage{algpseudocode}
\usepackage{float}

\onehalfspacing

\newcommand{\R}{\mathbb R}
\newcommand{\ind}{\mathbb I}
\newcommand{\norm}[1]{\left\lVert #1 \right\rVert}
\newcommand{\trans}{^{\mathsf T}}
\DeclareMathOperator{\logit}{logit}
\DeclareMathOperator{\Cat}{Categorical}
\DeclareMathOperator{\Dir}{Dirichlet}
\DeclareMathOperator{\IG}{IG}
\DeclareMathOperator{\IW}{IW}
\DeclareMathOperator{\tr}{tr}

\title{Joint Multilayered LSIRM with Latent-Position Gaussian Mixture Clustering\\
NIW Predictive Allocation and Jain--Neal Split--Merge Moves}
\date{Model Description}

\begin{document}
\maketitle

\section{Overview}

This document defines a joint multilayered LSIRM in which item clustering is
performed directly on item latent positions.  The previous PPCA measurement
block is removed.  In the modified model, each item latent position is assigned
a finite Gaussian mixture prior,
\[
    b_j^{(\ell(j))}\mid c_j=l,\mu_l,\Sigma_l
    \sim N_d(\mu_l,\Sigma_l),
    \qquad d=2,
\]
where $c_j$ is the item cluster label.  The mixture prior is imposed only on
the item latent positions.  It is not imposed directly on respondent ability or
item difficulty parameters.  Therefore, clustering regularizes the latent item
geometry while difficulty and ability remain governed by their own non-clustered
priors.

The model is fully joint: item cluster labels affect the posterior distribution
of item positions through the mixture prior, and item positions affect the
network/response likelihood through the LSIRM distance term.  Consequently,
cluster labels can indirectly affect respondent and item effects through the
posterior dependence induced by the likelihood, but no cluster-specific prior is
placed directly on those scalar effects.

\section{Indexing, dimensions, and observed data}

\begin{center}
\begin{tabular}{lll}
\toprule
Symbol & Meaning & Dimension/range \\
\midrule
$\ell$ & layer index & $\ell=1,\dots,L$ \\
$i$ & respondent/person index & $i=1,\dots,n_\ell$ \\
$j$ & local item index within layer $\ell$ & $j=1,\dots,P_\ell$ \\
$q$ & global item index across all layers & $q=1,\dots,P$, $P=\sum_{\ell=1}^L P_\ell$ \\
$\ell(q)$ & layer containing global item $q$ & $\{1,\dots,L\}$ \\
$j(q)$ & local item index of global item $q$ & $\{1,
\dots,P_{\ell(q)}\}$ \\
$d$ & latent space dimension & $d=2$ in the current implementation \\
$K^\star$ & truncation level / maximum number of mixture components & positive integer \\
$K_+$ & number of occupied mixture components & $K_+=|\{l:n_l>0\}|$ \\
\bottomrule
\end{tabular}
\end{center}

\paragraph{Common-respondent assumption.}
Throughout this document the typical multilayered setting with shared
respondents is assumed: $n_\ell=n$ for all $\ell$, and the respondent latent
position is shared across layers,
$a_i^{(\ell)}=a_i\in\R^d$ for $i=1,\dots,n$ and $\ell=1,\dots,L$.  The
layered notation $a_i^{(\ell)}$ is kept in the equations below for visual
symmetry with $b_j^{(\ell)}$ but should be read as the single $a_i$.  Two
consequences follow.  First, every $b_j^{(\ell)}$ is expressed in one
coordinate system anchored by the shared $a$'s; the only identifiability
ambiguity that remains is a global rigid motion of the whole configuration
$(a,b^{(1)},\dots,b^{(L)})$, not per-layer rigid motions.  Second, the
Metropolis--Hastings update of $a_i$ in \eqref{eq:a-mh-ratio} aggregates
likelihood contributions across all layers that contain respondent $i$, with
a single fixed unit-variance prior on $a_i$ (see \eqref{eq:a-prior}) rather
than $L$ layer-specific variances.  When the assumption fails (each layer
carries its own respondent set), the layered notation should be read
literally; rotational invariance then becomes per layer and stronger
alignment is required than the post-hoc step described in
Section~\ref{sec:layer-comparability}.

The observed data are denoted by
\[
    Y^{(\ell)}=\{y_{ij}^{(\ell)}:(i,j)\in\mathcal O_\ell\},
\]
where $\mathcal O_\ell$ is the set of observed respondent--item pairs in layer
$\ell$.  If there is no missingness, then
$\mathcal O_\ell=\{1,\dots,n_\ell\}\times\{1,\dots,P_\ell\}$.

\section{LSIRM likelihood}

For respondent $i$ and item $j$ in layer $\ell$, the model uses the following
parameters:
\begin{itemize}[leftmargin=1.5em]
    \item $a_i\in\R^d$: shared respondent latent position (same across layers).
    \item $b_j^{(\ell)}\in\R^d$: item latent position in layer $\ell$.
    \item $\alpha_i^{(\ell)}\in\R$: respondent effect in layer $\ell$.
    \item $\beta_j^{(\ell)}\in\R$: item difficulty for binary, continuous, and
    count layers.  In ordinal layers, $\beta_j^{(\ell)}$ is replaced by an
    item-specific threshold vector
    $\beta_{j,1}^{(\ell)}<\cdots<\beta_{j,K_\ell-1}^{(\ell)}$, where $K_\ell$
    is the number of categories.
    \item $\gamma_\ell>0$: layer-specific distance weight.
    \item $\sigma_0^2>0$: continuous-layer error scale.
    \item $\nu_2>0$: continuous-layer Student-$t$ degrees of freedom.
    \item $\kappa_j>0$: \emph{item-specific} count-layer overdispersion
    (negative binomial), one per count item $j$.
\end{itemize}
Let $\mathcal L_{\mathrm{ord}}\subseteq\{1,\dots,L\}$ index the ordinal
layers.

\paragraph{Implementation note.}
The general formulation above is for an arbitrary number of layers $L$.  The
current implementation is a four-layer instance ($L=4$), with one binary, one
continuous, one count, and one ordinal (LSGRM) layer; that is,
$\mathcal L_{\mathrm{ord}}=\{4\}$.  The notation below is kept layer-generic.

The layer-specific LSIRM linear predictor is
\begin{equation}
    \eta_{ij}^{(\ell)}
    =
    \alpha_i^{(\ell)}-\beta_j^{(\ell)}-\gamma_\ell\norm{a_i-b_j^{(\ell)}}_2,
    \qquad \ell\notin\mathcal L_{\mathrm{ord}},
    \label{eq:lsirm-linear-predictor}
\end{equation}
with the convention that for ordinal layers $\beta_j^{(\ell)}$ does not appear
in the linear predictor (the item thresholds enter the likelihood directly):
\begin{equation}
    \eta_{ij}^{(\ell)}
    =
    \alpha_i^{(\ell)}-\gamma_\ell\norm{a_i-b_j^{(\ell)}}_2,
    \qquad \ell\in\mathcal L_{\mathrm{ord}}.
    \label{eq:lsirm-linear-predictor-ord}
\end{equation}

\subsection{Per-layer observation models}

\paragraph{Binary (logistic regression).}
For a binary layer $\ell$,
\begin{equation}
    y_{ij}^{(\ell)}\mid \eta_{ij}^{(\ell)}
    \sim \mathrm{Bernoulli}(p_{ij}^{(\ell)}),
    \qquad
    \logit p_{ij}^{(\ell)}=\eta_{ij}^{(\ell)}.
    \label{eq:lik-binary}
\end{equation}

\paragraph{Continuous (robust Student-$t$ via normal--gamma scale mixture).}
For a continuous layer $\ell$,
\begin{align}
    y_{ij}^{(\ell)}\mid \eta_{ij}^{(\ell)},\lambda_{ij}^{(\ell)}
    &\sim N\!\left(\eta_{ij}^{(\ell)},\;\sigma_0^2/\lambda_{ij}^{(\ell)}\right),
    \notag\\
    \lambda_{ij}^{(\ell)}
    &\sim \mathrm{Gamma}(\nu_2/2,\;\nu_2/2).
    \label{eq:lik-continuous}
\end{align}
Marginalizing over $\lambda_{ij}^{(\ell)}$ gives
$y_{ij}^{(\ell)}\mid \eta_{ij}^{(\ell)}\sim t_{\nu_2}(\eta_{ij}^{(\ell)},\sigma_0^2)$,
the Student-$t$ regression observation model.

\paragraph{Count (negative binomial regression with item-specific overdispersion).}
For a count layer $\ell$, each item $j$ has its own overdispersion
parameter $\kappa_j>0$:
\begin{equation}
    y_{ij}^{(\ell)}\mid \eta_{ij}^{(\ell)},\kappa_j
    \sim \mathrm{NB}\!\left(r=\tfrac{1}{\kappa_j},\;\mu=\exp\eta_{ij}^{(\ell)}\right),
    \label{eq:lik-count}
\end{equation}
parameterized so that
$E[y_{ij}^{(\ell)}\mid\eta_{ij}^{(\ell)}]=\exp\eta_{ij}^{(\ell)}$ and
$\mathrm{Var}[y_{ij}^{(\ell)}\mid\eta_{ij}^{(\ell)}]=\mu+\kappa_j\mu^2$.
The Poisson regression limit is recovered as $\kappa_j\to 0$.  The
overdispersion is taken to be item-specific so that items with very different
dispersion patterns are not forced to share a single $\kappa$.  We assign
each $\kappa_j$ a log-normal prior,
\begin{equation}
    \log\kappa_j\sim N(\mu_{\log\kappa},\sigma_{\log\kappa}^2),
    \label{eq:kappa-prior}
\end{equation}
with $(\mu_{\log\kappa},\sigma_{\log\kappa}^2)$ fixed hyperparameters.

\paragraph{Ordinal (latent-space graded response model, LSGRM).}
We follow the LSGRM specification of De~Carolis, Kang, and Jeon (2025).  For
an ordinal layer $\ell\in\mathcal L_{\mathrm{ord}}$ with $K_\ell$ ordered
categories ($y_{ij}^{(\ell)}\in\{1,\dots,K_\ell\}$) and item thresholds
$\beta_{j,1}^{(\ell)}<\cdots<\beta_{j,K_\ell-1}^{(\ell)}$, the cumulative
response probabilities are
\begin{equation}
    \logit
    P\bigl(y_{ij}^{(\ell)}\geq k+1\bigm|\eta_{ij}^{(\ell)},\beta_{j,k}^{(\ell)}\bigr)
    =
    \eta_{ij}^{(\ell)}-\beta_{j,k}^{(\ell)}
    =
    \alpha_i^{(\ell)}-\beta_{j,k}^{(\ell)}-\gamma_\ell\norm{a_i-b_j^{(\ell)}}_2,
    \qquad k=1,\dots,K_\ell-1,
    \label{eq:lik-ordinal-cumulative}
\end{equation}
with the boundary conventions $P(y_{ij}^{(\ell)}\geq 1)=1$ and
$P(y_{ij}^{(\ell)}\geq K_\ell+1)=0$.  The category probabilities are obtained
by differencing
\begin{equation}
    P\bigl(y_{ij}^{(\ell)}=k\bigm|\cdot\bigr)
    =
    P\bigl(y_{ij}^{(\ell)}\geq k\bigm|\cdot\bigr)
    -P\bigl(y_{ij}^{(\ell)}\geq k+1\bigm|\cdot\bigr),
    \qquad k=1,\dots,K_\ell.
    \label{eq:lik-ordinal}
\end{equation}
The threshold sign and ordering convention here matches De~Carolis et al.\
(2025): a larger $\beta_{j,k}^{(\ell)}$ shifts more probability mass toward
lower categories, and the ordered support prevents label-flipping between
adjacent categories.

\subsection{Compact notation}

Each per-layer likelihood is summarized as
\begin{equation}
    y_{ij}^{(\ell)}\mid \eta_{ij}^{(\ell)},\xi_\ell
    \sim f_\ell\bigl\{y_{ij}^{(\ell)}\mid \eta_{ij}^{(\ell)},\xi_\ell\bigr\},
    \qquad (i,j)\in\mathcal O_\ell,
    \label{eq:generic-layer-likelihood}
\end{equation}
where the layer-specific nuisance parameter $\xi_\ell$ collects, for example,
$(\sigma_0^2,\lambda_{ij}^{(\ell)})$ in continuous layers,
$\{\kappa_j:j=1,\dots,P_\ell\}$ in count layers, and the threshold vector
$(\beta_{j,1}^{(\ell)},\dots,\beta_{j,K_\ell-1}^{(\ell)})$ in ordinal layers,
together with $\gamma_\ell$ in every layer.  The full LSIRM log-likelihood is
\begin{equation}
    \ell_{\mathrm{LSIRM}}
    =
    \sum_{\ell=1}^L
    \sum_{(i,j)\in\mathcal O_\ell}
    \log f_\ell\bigl\{y_{ij}^{(\ell)}\mid \eta_{ij}^{(\ell)},\xi_\ell\bigr\}.
    \label{eq:lsirm-loglik}
\end{equation}

\section{Latent-position Gaussian mixture prior for items}

\subsection{Global item position notation}

Define the global item latent position
\begin{equation}
    z_q
    =
    b_{j(q)}^{(\ell(q))}
    \in\R^d,
    \qquad q=1,\dots,P.
    \label{eq:global-item-position}
\end{equation}
The clustering model is placed directly on $z_q$.

\subsection{Finite Gaussian mixture prior}

For $q=1,\dots,P$ and $l=1,\dots,K^\star$,
\begin{align}
    z_q\mid c_q=l,\mu_l,\Sigma_l
    &\sim N_d(\mu_l,\Sigma_l),
    \label{eq:z-mixture}\\
    c_q\mid \rho
    &\sim \Cat(\rho_1,\dots,\rho_{K^\star}),\\
    \rho
    &\sim \Dir(e_0,\dots,e_0).
    \label{eq:rho-prior}
\end{align}

The mixture component parameters follow a Normal--Inverse-Wishart prior:
\begin{align}
    \Sigma_l
    &\sim \IW(\nu_0,S_0),
    \label{eq:sigma-niw-prior}\\
    \mu_l\mid \Sigma_l
    &\sim N_d\left(m_0,\frac{1}{\kappa_0}\Sigma_l\right),
    \qquad l=1,\dots,K^\star.
    \label{eq:mu-niw-prior}
\end{align}
Here $m_0\in\R^d$, $\kappa_0>0$, $\nu_0>d-1$, and
$S_0\in\R^{d\times d}$ is positive definite.

\paragraph{Data-driven cluster scale: Wishart hyperprior on $S_0$.}
Rather than fixing $S_0$ at a calibrated value, we put a Wishart hyperprior on it,
\begin{equation}
    S_0\sim W(\nu_{S_0},\Lambda_{S_0}),
    \qquad \nu_{S_0}>d-1,\;\Lambda_{S_0}\in\R^{d\times d}\;\text{positive definite},
    \label{eq:S0-hyperprior}
\end{equation}
so that the cluster-covariance scale is learned from data instead of being a
manually tuned hyperparameter.  The Wishart is conjugate for the IW scale
matrix: combining \eqref{eq:S0-hyperprior} with
$\Sigma_l\mid S_0\sim\IW(\nu_0,S_0)$ for $l=1,\dots,K^\star$ gives the
closed-form full conditional
\begin{equation}
    S_0\mid\{\Sigma_l\}_{l=1}^{K^\star}
    \sim
    W\!\left(
    \nu_{S_0}+K^\star\nu_0,\;
    \bigl(\Lambda_{S_0}^{-1}+\sum_{l=1}^{K^\star}\Sigma_l^{-1}\bigr)^{-1}
    \right),
    \label{eq:S0-update}
\end{equation}
which we sample by ordinary Gibbs once per MCMC sweep (after the
$(\mu_l,\Sigma_l)$ draws in \eqref{eq:sigma-posterior}--\eqref{eq:mu-posterior}).
Defaults are weakly informative and anchored at a user-provided initial
$S_0^{(0)}$:
$\nu_{S_0}=d+2$ and $\Lambda_{S_0}=S_0^{(0)}/\nu_{S_0}$, yielding
$E[S_0]=S_0^{(0)}$.  Setting $\nu_{S_0}\to\infty$ recovers the fixed-$S_0$
behaviour as a limiting case.

\paragraph{Important modeling restriction.}
The cluster label $c_q$ is allowed to affect $z_q=b_{j(q)}^{(\ell(q))}$ only.
There is no cluster-specific prior of the form
$\alpha_i^{(\ell)}\mid c_q$, $\beta_j^{(\ell)}\mid c_q$, or
$\xi_\ell\mid c_q$.  Thus, clusters represent latent item geometry rather than
marginal difficulty, ability, or layer-specific nuisance effects.

\section{Priors for non-clustering LSIRM parameters}

\paragraph{Respondent latent position (fixed unit variance for identifiability).}
The shared respondent latent position is given a standard normal prior with
fixed unit variance,
\begin{equation}
    a_i \sim N_d(0,I_d),\qquad i=1,\dots,n.
    \label{eq:a-prior}
\end{equation}
The variance is fixed rather than sampled because the latent space has no
intrinsic scale: with $a_i$ and $b_j^{(\ell)}$ both scalable by an arbitrary
constant absorbed into $\gamma_\ell$, only the \emph{ratio} of position scale
to $\gamma_\ell$ is identified.  Following the LSGRM convention of
De~Carolis, Kang, and Jeon (2025, \S2.4), we anchor the latent-space scale by
fixing the respondent-position prior variance to $1$, which leaves
$\gamma_\ell$ as the free scale parameter to be estimated from data.

\paragraph{Respondent effect.}
The respondent effect is layer-specific with its own variance hyperparameter,
\begin{equation}
    \alpha_i^{(\ell)}\mid \sigma_{\alpha,\ell}^2
    \sim N(0,\sigma_{\alpha,\ell}^2),
    \qquad i=1,\dots,n,\;\ell=1,\dots,L,
    \label{eq:alpha-prior}
\end{equation}
where $\sigma_{\alpha,\ell}^2$ is sampled from a layer-specific inverse-gamma
hyperprior:
\begin{equation}
    \sigma_{\alpha,\ell}^2 \sim \IG(A_{\alpha,\ell},B_{\alpha,\ell}),
    \qquad \ell=1,\dots,L.
    \label{eq:sigma-alpha-prior}
\end{equation}
Allowing the variance to differ across layers reflects the fact that the
overall scale of the respondent effect can differ between, e.g., binary and
count layers.

\paragraph{Item difficulty $\beta_j^{(\ell)}$ (fixed-effect, fixed-scale prior).}
For non-ordinal layers, the item difficulty is treated as a fixed effect:
\begin{equation}
    \beta_j^{(\ell)}\sim N(0,\sigma_{\beta,\ell}^2),
    \qquad j=1,\dots,P_\ell,\;\ell\notin\mathcal L_{\mathrm{ord}},
    \label{eq:beta-prior}
\end{equation}
where $\sigma_{\beta,\ell}^2$ is a \emph{fixed} hyperparameter (not sampled).
No shrinkage hierarchy is imposed, so each item's difficulty is recovered
directly from its own data.  The current implementation uses $\sigma_{\beta,\ell}=2$
for every non-ordinal layer (i.e.\ variance $4$).

\paragraph{Ordinal thresholds (independent normals subject to ordering).}
For ordinal layers $\ell\in\mathcal L_{\mathrm{ord}}$, each threshold is
assigned an independent normal prior with the same fixed scale,
\begin{equation}
    \beta_{j,k}^{(\ell)}\sim N(0,\sigma_{\beta,\ell}^2),
    \qquad j=1,\dots,P_\ell,\;k=1,\dots,K_\ell-1,
    \label{eq:beta-thr-prior}
\end{equation}
restricted to the ordered support
$\beta_{j,1}^{(\ell)}<\cdots<\beta_{j,K_\ell-1}^{(\ell)}$.  The scale
$\sigma_{\beta,\ell}^2$ is again fixed (not sampled); the implementation uses
$\sigma_{\beta,\ell}=2$ for the ordinal layer as well.  Independent normals
with an order constraint match the LSGRM update used in the implementation,
which proposes one threshold at a time (rejecting any move that violates the
ordering); a joint multivariate-normal prior with the same diagonal scale
would give an equivalent posterior.

The item position $b_j^{(\ell)}$ does not receive an additional independent
Gaussian prior; its prior is the latent-position mixture prior in
\eqref{eq:z-mixture}--\eqref{eq:mu-niw-prior}.

The remaining layer-specific nuisance parameters --- $\gamma_\ell$ (distance
weight, all layers), $\sigma_0^2$ (continuous-layer error scale), and
$\{\kappa_j\}$ (count-layer item-specific overdispersion, see
\eqref{eq:kappa-prior}) --- are equipped with independent log-normal or
inverse-gamma priors as appropriate; the precise form follows the V6/V8
specification and does not interact with the clustering layer.

\section{Rotation and reflection invariance}
\label{sec:layer-comparability}

Under the common-respondent assumption above, the entire configuration
$(a,b^{(1)},\dots,b^{(L)})$ shares one coordinate system, so the only
identifiability ambiguity in the latent space is a single global orthogonal
transformation $Q\in O(d)$ together with a single global translation
$t\in\R^d$:
\begin{equation}
    a_i\mapsto Qa_i+t,
    \qquad
    b_j^{(\ell)}\mapsto Qb_j^{(\ell)}+t,
    \qquad
    \mu_l\mapsto Q\mu_l+t,
    \qquad
    \Sigma_l\mapsto Q\Sigma_l Q\trans.
    \label{eq:rigid-motion}
\end{equation}
The LSIRM linear predictor in
\eqref{eq:lsirm-linear-predictor}--\eqref{eq:lsirm-linear-predictor-ord}
depends on positions only through $\norm{a_i-b_j^{(\ell)}}_2$, which is
invariant under \eqref{eq:rigid-motion}.  With $m_0=0$ and $S_0=s_0 I_d$, the
NIW prior on $(\mu_l,\Sigma_l)$ in
\eqref{eq:sigma-niw-prior}--\eqref{eq:mu-niw-prior} is also invariant in
distribution under $Q\in O(d)$.  Therefore the joint posterior on the
clustering targets---$c$, $K_+$, the cluster sizes $n_l$, and the posterior
similarity matrix $C_{qr}=P(c_q=c_r\mid Y)$---is invariant under
\eqref{eq:rigid-motion}.

\paragraph{Consequence for the sampler.}
Within-MCMC orientation alignment of the latent space is \emph{not} required
for clustering correctness.  The chain may drift in orientation from one
iteration to the next without affecting any cluster-label-based posterior
summary.  A post-hoc Procrustes alignment is recommended only when the
\emph{visualization} of $a_i$ or $b_j^{(\ell)}$ must be compared to a target
configuration (e.g.\ true positions in a simulation, or a reference set in
applied work).  Such an alignment is a deterministic post-processing step
applied to stored posterior draws and does not enter the MCMC update for
$c$, $\mu_l$, $\Sigma_l$, or $\rho$.

\section{Full joint posterior}

Let
\begin{align*}
    \Theta_{\mathrm{LSIRM}}
    &=
    \{a_i,\alpha_i^{(\ell)},b_j^{(\ell)},\beta_j^{(\ell)},\xi_\ell,
    \sigma_{\alpha,\ell}^2:
    i=1,\dots,n,\ \ell=1,\dots,L\},\\
    \Theta_{\mathrm{mix}}
    &=
    \{c_q,\rho,\mu_l,\Sigma_l:q=1,\dots,P,\ l=1,\dots,K^\star\}.
\end{align*}
The full posterior is proportional to
\begin{align}
&p(\Theta_{\mathrm{LSIRM}},\Theta_{\mathrm{mix}}\mid Y)
\notag\\
&\quad \propto
\prod_{\ell=1}^L
\prod_{(i,j)\in\mathcal O_\ell}
    f_\ell\{y_{ij}^{(\ell)}\mid \eta_{ij}^{(\ell)},\xi_\ell\}
\notag\\
&\qquad\times
\prod_{q=1}^P
    N_d(z_q;\mu_{c_q},\Sigma_{c_q})\rho_{c_q}
\notag\\
&\qquad\times
\prod_{l=1}^{K^\star}
    p(\mu_l\mid\Sigma_l)p(\Sigma_l)
    \ p(\rho)
\notag\\
&\qquad\times
\prod_{i=1}^{n}N_d(a_i;0,I_d)
\prod_{\ell=1}^L
\left[
\prod_{i=1}^{n}p(\alpha_i^{(\ell)}\mid\sigma_{\alpha,\ell}^2)
\;p_\beta^{(\ell)}\bigl(\beta^{(\ell)}\bigr)
\right]
\notag\\
&\qquad\times
\prod_{\ell=1}^L
    p(\sigma_{\alpha,\ell}^2)p(\xi_\ell),
    \label{eq:full-posterior}
\end{align}
where $z_q=b_{j(q)}^{(\ell(q))}$ and the layer-specific item-difficulty prior
factor is
\[
    p_\beta^{(\ell)}\bigl(\beta^{(\ell)}\bigr)
    =
    \begin{cases}
    \displaystyle\prod_{j=1}^{P_\ell} N\!\bigl(\beta_j^{(\ell)};0,\sigma_{\beta,\ell}^2\bigr),
    & \ell\notin\mathcal L_{\mathrm{ord}},\\[1em]
    \displaystyle\prod_{j=1}^{P_\ell}
    \left[
    \prod_{k=1}^{K_\ell-1}
    N\!\bigl(\beta_{j,k}^{(\ell)};0,\sigma_{\beta,\ell}^2\bigr)
    \right]
    \ind\bigl\{\beta_{j,1}^{(\ell)}<\cdots<\beta_{j,K_\ell-1}^{(\ell)}\bigr\},
    & \ell\in\mathcal L_{\mathrm{ord}}.
    \end{cases}
\]
The respondent-position prior contributes $N_d(a_i;0,I_d)$ with fixed unit
variance.  The item-difficulty scales $\sigma_{\beta,\ell}^2$ are fixed
hyperparameters and enter as constants in the posterior; the implementation
sets $\sigma_{\beta,\ell}=2$ throughout.

\section{NIW posterior quantities}

For a generic set of global item indices $A\subseteq\{1,\dots,P\}$, define
\begin{align}
    n_A &= |A|,\\
    \bar z_A &= \frac{1}{n_A}\sum_{q\in A}z_q,
    \qquad n_A>0,\\
    Q_A &= \sum_{q\in A}(z_q-\bar z_A)(z_q-\bar z_A)\trans,
    \qquad n_A>0.
\end{align}
For $n_A=0$, use prior quantities directly and set $Q_A=0$.

The NIW posterior parameters for the items in $A$ are
\begin{align}
    \kappa_A &= \kappa_0+n_A,\\
    \nu_A &= \nu_0+n_A,\\
    m_A &=
    \begin{cases}
    \dfrac{\kappa_0m_0+n_A\bar z_A}{\kappa_0+n_A}, & n_A>0,\\[1em]
    m_0, & n_A=0,
    \end{cases}\\
    S_A &=
    \begin{cases}
    S_0+Q_A+
    \dfrac{\kappa_0 n_A}{\kappa_0+n_A}
    (\bar z_A-m_0)(\bar z_A-m_0)\trans,
    & n_A>0,\\[1em]
    S_0, & n_A=0.
    \end{cases}
    \label{eq:niw-posterior-params}
\end{align}

For component $l$, let
\[
    \mathcal J_l=\{q:c_q=l\},
    \qquad n_l=|\mathcal J_l|.
\]
Then
\[
    (\kappa_l,\nu_l,m_l,S_l)
    =
    (\kappa_{\mathcal J_l},\nu_{\mathcal J_l},m_{\mathcal J_l},S_{\mathcal J_l}).
\]
For leave-one-out calculations,
\[
    \mathcal J_{l,-q}=\{r:r\neq q,\ c_r=l\}.
\]

\section{NIW posterior predictive allocation update}

Integrating out $(\mu_l,\Sigma_l)$ for label allocation gives the Student-$t$
posterior predictive density
\begin{equation}
    z_q\mid z_{\mathcal J_{l,-q}}
    \sim
    t_{\nu_{l,-q}-d+1}
    \left(
    m_{l,-q},
    \frac{\kappa_{l,-q}+1}
    {\kappa_{l,-q}(\nu_{l,-q}-d+1)}S_{l,-q}
    \right).
    \label{eq:niw-predictive}
\end{equation}
Therefore the collapsed single-site label update is
\begin{equation}
    P(c_q=l\mid c_{-q},z)
    \propto
    (n_{l,-q}+e_0)
    t_{\nu_{l,-q}-d+1}
    \left(
    z_q;
    m_{l,-q},
    \frac{\kappa_{l,-q}+1}
    {\kappa_{l,-q}(\nu_{l,-q}-d+1)}S_{l,-q}
    \right),
    \label{eq:c-update}
\end{equation}
for $l=1,\dots,K^\star$.  This is a Gibbs update, so the acceptance
probability is one after sampling from the normalized probabilities in
\eqref{eq:c-update}.

\section{Posterior draw of mixture component parameters}

After updating the labels, draw component parameters from their NIW posterior:
\begin{align}
    \Sigma_l\mid z,c
    &\sim \IW(\nu_l,S_l),
    \label{eq:sigma-posterior}\\
    \mu_l\mid \Sigma_l,z,c
    &\sim N_d\left(m_l,\frac{1}{\kappa_l}\Sigma_l\right),
    \qquad l=1,\dots,K^\star.
    \label{eq:mu-posterior}
\end{align}
For empty components, $n_l=0$, and this reduces to a draw from the NIW prior.

After all $\Sigma_l$ have been refreshed, sample the cluster-scale
hyperparameter from its Wishart full conditional in
\eqref{eq:S0-update}:
\begin{equation}
    S_0\mid\{\Sigma_l\}
    \sim
    W\!\left(
    \nu_{S_0}+K^\star\nu_0,\;
    \bigl(\Lambda_{S_0}^{-1}+\textstyle\sum_l\Sigma_l^{-1}\bigr)^{-1}
    \right).
    \label{eq:S0-gibbs}
\end{equation}
The cached log-determinant $\log|S_0|$ used in the NIW marginal likelihood
\eqref{eq:niw-marginal} for split--merge moves must be refreshed after every
$S_0$ draw.

The mixture weights may be sampled for monitoring or storage:
\begin{equation}
    \rho\mid c
    \sim
    \Dir(e_0+n_1,\dots,e_0+n_{K^\star}).
    \label{eq:rho-posterior}
\end{equation}
The draw of $\rho$ is not needed for the collapsed label update.

\section{Collapsed partition target for split--merge moves}

Integrating out $\rho$ and all component parameters gives the collapsed
partition target
\begin{equation}
    \pi(c\mid z)
    \propto
    \prod_{l=1}^{K^\star}
    \frac{\Gamma(n_l+e_0)}{\Gamma(e_0)}
    m(z_{\mathcal J_l}),
    \label{eq:collapsed-partition-target}
\end{equation}
where
\begin{equation}
    m(z_A)
    =
    \int
    \prod_{q\in A}N_d(z_q;\mu,\Sigma)
    p(\mu,\Sigma)\,d\mu\,d\Sigma
\end{equation}
is the NIW marginal likelihood.

The log NIW marginal likelihood is
\begin{align}
    \log m(z_A)
    &=-\frac{n_A d}{2}\log\pi
    +\frac{d}{2}\{\log\kappa_0-\log\kappa_A\}
    +\frac{\nu_0}{2}\log|S_0|
    -\frac{\nu_A}{2}\log|S_A|
    \notag\\
    &\qquad
    +\log\Gamma_d\left(\frac{\nu_A}{2}\right)
    -\log\Gamma_d\left(\frac{\nu_0}{2}\right),
    \label{eq:niw-marginal}
\end{align}
where
\begin{equation}
    \log\Gamma_d(a)
    =
    \frac{d(d-1)}{4}\log\pi
    +
    \sum_{h=1}^d\log\Gamma\left(a+\frac{1-h}{2}\right).
\end{equation}
Thus,
\begin{equation}
    \log\pi(c\mid z)
    =
    \sum_{l=1}^{K^\star}
    \left[
    \log\Gamma(n_l+e_0)-\log\Gamma(e_0)
    +\log m(z_{\mathcal J_l})
    \right]
    +\mathrm{constant}.
    \label{eq:log-collapsed-partition-target}
\end{equation}

\section{Metropolis--Hastings updates for LSIRM parameters}

All LSIRM parameters with non-conjugate full conditionals can be updated by
random-walk Metropolis--Hastings.  The formulas below assume symmetric
Gaussian random-walk proposals.  If a non-symmetric proposal is used, the
corresponding proposal-density ratio must be added to the log acceptance ratio.

\subsection{Respondent effect $\alpha_i^{(\ell)}$}

Propose
\[
    \alpha_i^{(\ell)\prime}
    =
    \alpha_i^{(\ell)}+\varepsilon_\alpha,
    \qquad
    \varepsilon_\alpha\sim N(0,s_\alpha^2).
\]
Let $\eta_{ij}^{(\ell)\prime}$ be the linear predictor after replacing
$\alpha_i^{(\ell)}$ by $\alpha_i^{(\ell)\prime}$.  The log acceptance ratio is
\begin{align}
    \log R_\alpha
    &=
    \sum_{j:(i,j)\in\mathcal O_\ell}
    \left[
    \log f_\ell\{y_{ij}^{(\ell)}\mid \eta_{ij}^{(\ell)\prime},\xi_\ell\}
    -
    \log f_\ell\{y_{ij}^{(\ell)}\mid \eta_{ij}^{(\ell)},\xi_\ell\}
    \right]
    \notag\\
    &\qquad
    -
    \frac{1}{2\sigma_{\alpha,\ell}^2}
    \left\{(\alpha_i^{(\ell)\prime})^2-(\alpha_i^{(\ell)})^2\right\}.
    \label{eq:alpha-mh-ratio}
\end{align}
The proposal is accepted with probability $\min\{1,\exp(\log R_\alpha)\}$.

\subsection{Item difficulty $\beta_j^{(\ell)}$}

For non-ordinal layers $\ell\notin\mathcal L_{\mathrm{ord}}$, propose
\[
    \beta_j^{(\ell)\prime}
    =
    \beta_j^{(\ell)}+\varepsilon_\beta,
    \qquad
    \varepsilon_\beta\sim N(0,s_\beta^2).
\]
With the fixed-scale prior \eqref{eq:beta-prior}, the log acceptance ratio is
\begin{align}
    \log R_\beta
    &=
    \sum_{i:(i,j)\in\mathcal O_\ell}
    \left[
    \log f_\ell\{y_{ij}^{(\ell)}\mid \eta_{ij}^{(\ell)\prime},\xi_\ell\}
    -
    \log f_\ell\{y_{ij}^{(\ell)}\mid \eta_{ij}^{(\ell)},\xi_\ell\}
    \right]
    \notag\\
    &\qquad
    -\frac{1}{2\sigma_{\beta,\ell}^2}
    \left\{(\beta_j^{(\ell)\prime})^2-(\beta_j^{(\ell)})^2\right\}.
    \label{eq:beta-mh-ratio}
\end{align}

For ordinal layers $\ell\in\mathcal L_{\mathrm{ord}}$, propose each threshold
$\beta_{j,k}^{(\ell)}$ ($k=1,\dots,K_\ell-1$) by a Gaussian random walk with
its own scale $s_{\beta,k}^2$.  The proposed move is rejected outright if the
ordering $\beta_{j,1}^{(\ell)}<\cdots<\beta_{j,K_\ell-1}^{(\ell)}$ is
violated.  Otherwise, with the fixed-scale independent-normal prior
\eqref{eq:beta-thr-prior}, the log acceptance ratio is
\begin{align}
    \log R_{\beta,k}
    &=
    \sum_{i:(i,j)\in\mathcal O_\ell}
    \left[
    \log f_\ell\{y_{ij}^{(\ell)}\mid \eta_{ij}^{(\ell)},\beta_{j,\cdot}^{(\ell)\prime}\}
    -
    \log f_\ell\{y_{ij}^{(\ell)}\mid \eta_{ij}^{(\ell)},\beta_{j,\cdot}^{(\ell)}\}
    \right]
    \notag\\
    &\qquad
    -\frac{1}{2\sigma_{\beta,\ell}^2}
    \left\{(\beta_{j,k}^{(\ell)\prime})^2-(\beta_{j,k}^{(\ell)})^2\right\},
    \label{eq:beta-thr-mh-ratio}
\end{align}
where the layer-$\ell$ ordinal likelihood
\eqref{eq:lik-ordinal-cumulative}--\eqref{eq:lik-ordinal} depends on the full
threshold vector $\beta_{j,1:K_\ell-1}^{(\ell)}$ for item $j$ rather than on
$\eta_{ij}^{(\ell)}$ alone.

\subsection{Respondent latent position $a_i$}

Because $a_i$ is shared across layers, propose a single update
\[
    a_i'=a_i+\varepsilon_a,\qquad \varepsilon_a\sim N_d(0,s_a^2 I_d),
\]
and aggregate likelihood contributions from every layer that contains
respondent $i$:
\begin{align}
    \log R_a
    &=
    \sum_{\ell=1}^L
    \sum_{j:(i,j)\in\mathcal O_\ell}
    \left[
    \log f_\ell\{y_{ij}^{(\ell)}\mid \eta_{ij}^{(\ell)\prime},\xi_\ell\}
    -
    \log f_\ell\{y_{ij}^{(\ell)}\mid \eta_{ij}^{(\ell)},\xi_\ell\}
    \right]
    \notag\\
    &\qquad
    -\tfrac{1}{2}\left\{\norm{a_i'}_2^2-\norm{a_i}_2^2\right\},
    \label{eq:a-mh-ratio}
\end{align}
where $\eta_{ij}^{(\ell)\prime}$ is the layer-$\ell$ linear predictor
evaluated at $a_i'$.  No $\sigma_a^2$ appears because the prior variance is
fixed at $1$ for identifiability \eqref{eq:a-prior}.

\subsection{Item latent position $b_j^{(\ell)}$}

Let $q=q(\ell,j)$ be the global item index corresponding to local item $j$ in
layer $\ell$, and let $c_q$ be its current cluster label.  Propose
\[
    b_j^{(\ell)\prime}
    =
    b_j^{(\ell)}+\varepsilon_b,
    \qquad
    \varepsilon_b\sim N_d(0,s_b^2 I_d).
\]
The log acceptance ratio includes both the LSIRM likelihood contribution and
the Gaussian mixture prior contribution:
\begin{align}
    \log R_b
    &=
    \sum_{i:(i,j)\in\mathcal O_\ell}
    \left[
    \log f_\ell\{y_{ij}^{(\ell)}\mid \eta_{ij}^{(\ell)\prime},\xi_\ell\}
    -
    \log f_\ell\{y_{ij}^{(\ell)}\mid \eta_{ij}^{(\ell)},\xi_\ell\}
    \right]
    \notag\\
    &\qquad
    +
    \log N_d\{b_j^{(\ell)\prime};\mu_{c_q},\Sigma_{c_q}\}
    -
    \log N_d\{b_j^{(\ell)};\mu_{c_q},\Sigma_{c_q}\}.
    \label{eq:b-mh-ratio}
\end{align}
This is the main difference from a post-hoc clustering approach: the cluster
assignment $c_q$ enters the posterior update of $b_j^{(\ell)}$ through the
mixture prior term in \eqref{eq:b-mh-ratio}.

\paragraph{Sampler-level variance inflation in the $b$ update (optional).}
The full coupling in \eqref{eq:b-mh-ratio} ties two roles to the same cluster
covariance $\Sigma_{c_q}$: (i) defining cluster identity (via the c, $\mu$,
$\Sigma$ updates) and (ii) anchoring $b_j^{(\ell)}$'s LSIRM geometry.
Tight $\Sigma_{c_q}$ (small posterior cluster covariance) gives a strong $b$
pull toward $\mu_{c_q}$, which compresses within-cluster pairwise distances
and elongates between-cluster distances --- visible as an S-curve distortion
in the distance scatter.

A sampler-level remedy is to inflate the $b$ MH prior covariance by a fixed
factor $\alpha\geq 1$ \emph{only inside the $b$ update}:
\begin{equation}
    \log R_b
    =
    \sum_{i:(i,j)\in\mathcal O_\ell}
    \bigl[\log f_\ell(\cdot\mid \eta_{ij}^{(\ell)\prime})
        - \log f_\ell(\cdot\mid \eta_{ij}^{(\ell)})\bigr]
    -
    \frac{1}{2\alpha}
    \bigl[(b_j^{(\ell)\prime}-\mu_{c_q})\trans \Sigma_{c_q}^{-1}(b_j^{(\ell)\prime}-\mu_{c_q})
    -(b_j^{(\ell)}-\mu_{c_q})\trans \Sigma_{c_q}^{-1}(b_j^{(\ell)}-\mu_{c_q})\bigr],
    \label{eq:b-mh-ratio-inflated}
\end{equation}
equivalent to using prior $b_j^{(\ell)}\sim N_d(\mu_{c_q},\alpha\Sigma_{c_q})$
in the MH ratio.  All other updates ($c$, $\mu_l$, $\Sigma_l$, $S_0$) use
the original $\Sigma_{c_q}$ ($\alpha=1$), so cluster identity, NIW posterior
draws, and split--merge marginal likelihoods are unchanged.

Limits:
\begin{itemize}[leftmargin=1.5em]
    \item $\alpha=1$: paper-defined fully joint posterior \eqref{eq:b-mh-ratio}.
    \item $\alpha=+\infty$: no $b$ pull, equivalent to a one-way coupling
    where clustering is informed by $b$ but $b$ is determined by LSIRM alone.
    \item $\alpha\in(1,\infty)$: tunable interpolation that loosens the
    geometry anchor while preserving cluster identity.
\end{itemize}

This is a sampler-level approximation, not a change to the joint model in
\eqref{eq:full-posterior}.  Formally, $\alpha>1$ corresponds to running the
$b$-block of a tempered ($1/\alpha$-power) version of the mixture prior
contribution while keeping the rest of the chain at the unit temperature.
Empirically it is useful when the LSIRM geometry must be reported alongside
cluster labels, since the within-cluster compression caused by $\alpha=1$
distorts pairwise-distance recovery without improving partition recovery.

\paragraph{Marginalized (collapsed-label) $b$ update.}
A more principled alternative to the tempering in
\eqref{eq:b-mh-ratio-inflated} is to integrate the cluster label $c_q$, the
weights $\rho$, and the component parameters $(\mu_l,\Sigma_l)$ out of the
prior factor of \eqref{eq:b-mh-ratio} entirely.  Since the NIW posterior
predictive in \eqref{eq:niw-predictive} and the collapsed allocation weights
in \eqref{eq:c-update} are already maintained by the chain, the leave-one-out
mixture predictive of $z_q$ given the remaining items and labels is
\begin{equation}
    p(z_q\mid z_{-q},c_{-q})
    \propto
    \sum_{l=1}^{K^\star}
    (n_{l,-q}+e_0)\;
    t_{\nu_{l,-q}-d+1}\!\left(
    z_q;\,m_{l,-q},\,
    \frac{\kappa_{l,-q}+1}{\kappa_{l,-q}(\nu_{l,-q}-d+1)}\,S_{l,-q}
    \right).
    \label{eq:b-marginal-prior}
\end{equation}
Define
\begin{equation}
    A(z;c_{-q},z_{-q})
    =
    \log\!\sum_{l=1}^{K^\star}
    \exp\!\Bigl[
        \log(n_{l,-q}+e_0)
        +\log t_{\nu_{l,-q}-d+1}\!\bigl(
            z;\,m_{l,-q},\,
            \tfrac{\kappa_{l,-q}+1}{\kappa_{l,-q}(\nu_{l,-q}-d+1)}\,S_{l,-q}
        \bigr)
    \Bigr],
    \label{eq:b-marginal-A}
\end{equation}
evaluated by log-sum-exp for numerical stability.  With the same symmetric
Gaussian random walk $b_j^{(\ell)\prime}=b_j^{(\ell)}+\varepsilon_b$,
$\varepsilon_b\sim N_d(0,s_b^2 I_d)$, the log MH acceptance ratio replaces
the single per-cluster Gaussian prior difference in \eqref{eq:b-mh-ratio} by
a log-sum-exp difference of the leave-one-out mixture predictive,
\begin{equation}
    \log R_b^{\mathrm{coll}}
    =
    \sum_{i:(i,j)\in\mathcal O_\ell}
    \!\bigl[\log f_\ell(\cdot\mid \eta_{ij}^{(\ell)\prime})
        - \log f_\ell(\cdot\mid \eta_{ij}^{(\ell)})\bigr]
    +
    A(z_q';c_{-q},z_{-q})
    -
    A(z_q;c_{-q},z_{-q}),
    \label{eq:b-mh-ratio-marginal}
\end{equation}
where $z_q=b_j^{(\ell)}$ and $z_q'=b_j^{(\ell)\prime}$.  The leave-one-out
NIW summaries $(\kappa_{l,-q},\nu_{l,-q},m_{l,-q},S_{l,-q})$ depend only on
$(z_{-q},c_{-q})$ and are therefore computed once per item update and
reused for both endpoints of the proposal.  Empty components ($n_{l,-q}=0$)
contribute the NIW prior predictive at $(\kappa_0,\nu_0,m_0,S_0)$ with
weight $e_0$.  For $l\neq c_q$ the leave-one-out summaries equal the full
NIW posterior summaries, so only the cluster $l=c_q$ requires actual
removal of $z_q$ from its sufficient statistics.

\paragraph{Validity as a partially collapsed Gibbs move.}
The update \eqref{eq:b-mh-ratio-marginal} integrates $(c_q,\rho,\mu,\Sigma)$
out of the $z_q$ block while keeping $c_{-q}$ fixed, so its target is the
marginal full conditional of $z_q$ under the \emph{same} joint posterior
\eqref{eq:full-posterior} as $\alpha=1$ in \eqref{eq:b-mh-ratio}.  By the
partially collapsed Gibbs construction of van~Dyk and Park (2008), the chain
preserves the same stationary distribution provided every marginalized
variable is refreshed from its conditional before being used downstream.
The sweep order in Algorithms~1--2 already satisfies this: after the $b$
update we resample $c$ via \eqref{eq:c-update}, then $(\Sigma_l,\mu_l)$ via
\eqref{eq:sigma-posterior}--\eqref{eq:mu-posterior}, then $S_0$ via
\eqref{eq:S0-gibbs}, and optionally $\rho$ via \eqref{eq:rho-posterior}.

\paragraph{Comparison with $\alpha$-inflation.}
Whereas \eqref{eq:b-mh-ratio-inflated} targets a tempered
($1/\alpha$-power) version of the mixture prior contribution, the
marginalized update \eqref{eq:b-mh-ratio-marginal} targets the original
$\alpha=1$ posterior exactly.  The geometry-compression failure mode is
mitigated structurally rather than by a tuning parameter: $b_j^{(\ell)}$ is
no longer anchored to a single $\mu_{c_q}$ through a tight $\Sigma_{c_q}$
but instead feels the full mixture predictive, whose tails are inflated by
the Student-$t$ shape and by the between-component support.  The additional
cost is $O(K^\star)$ Student-$t$ evaluations per $b$ MH proposal in place
of the single Gaussian density in \eqref{eq:b-mh-ratio}, which for
$K^\star\sim 10$ and $P\sim 30$ is a negligible fraction of the sweep cost.
The three options for the $b$ update --- conditional \eqref{eq:b-mh-ratio},
$\alpha$-inflated \eqref{eq:b-mh-ratio-inflated}, and marginalized
\eqref{eq:b-mh-ratio-marginal} --- are mutually exclusive and selected by
a sampler-level flag.

\section{Conjugate updates for variance parameters}

The only random-effect variance that is sampled is the layer-specific
respondent-effect variance $\sigma_{\alpha,\ell}^2$, with full conditional
\begin{equation}
    \sigma_{\alpha,\ell}^2\mid \alpha^{(\ell)}
    \sim
    \IG\left(
    A_{\alpha,\ell}+\frac{n}{2},\;
    B_{\alpha,\ell}+\frac{1}{2}\sum_{i=1}^{n}(\alpha_i^{(\ell)})^2
    \right),
    \qquad \ell=1,\dots,L.
    \label{eq:sigma-alpha-update}
\end{equation}
The respondent-position prior variance is fixed at $1$ \eqref{eq:a-prior} for
latent-space identifiability and is not sampled.  The item-difficulty scales
$\sigma_{\beta,\ell}^2$ in \eqref{eq:beta-prior}--\eqref{eq:beta-thr-prior}
are also fixed hyperparameters.

\section{Jain--Neal restricted Gibbs split--merge moves}

Split--merge moves are Metropolis--Hastings updates on $c$ conditional on the
current item positions $z=(z_1,\dots,z_P)$.  The target is the collapsed
partition density in \eqref{eq:log-collapsed-partition-target}.

\subsection{Anchor selection}

Select two distinct global item indices $u$ and $v$ uniformly from
$\{1,\dots,P\}$.  If $c_u=c_v$, attempt a split move.  If $c_u\neq c_v$,
attempt a merge move.  The anchor selection probability is symmetric and
cancels in the Metropolis--Hastings ratio.

\subsection{Split proposal}

Suppose $c_u=c_v=a$.  Let
\[
    A=\{q:c_q=a\}
\]
be the cluster to be split.  Choose an empty component label
$b\in\{l:n_l=0\}$ uniformly.  If no empty component exists, skip the split.

Initialize two temporary clusters by forcing
\[
    c_u'=a,
    \qquad
    c_v'=b.
\]
For each $q\in A\setminus\{u,v\}$, assign $q$ to either $a$ or $b$ by a
restricted NIW-collapsed Gibbs scan:
\begin{align}
    w_q(a)&\propto
    (n_{a,-q}^{\mathrm{temp}}+e_0)
    p(z_q\mid z_{a,-q}^{\mathrm{temp}}),\\
    w_q(b)&\propto
    (n_{b,-q}^{\mathrm{temp}}+e_0)
    p(z_q\mid z_{b,-q}^{\mathrm{temp}}),
\end{align}
where the predictive density is the NIW Student-$t$ density in
\eqref{eq:niw-predictive}.  Let $p_q^{\mathrm{fwd}}$ be the normalized
probability of the assignment actually selected during the forward restricted
scan.  Then
\begin{equation}
    \log q(c'\mid c)
    =
    -\log E(c)
    +
    \sum_{q\in A\setminus\{u,v\}}\log p_q^{\mathrm{fwd}},
    \label{eq:split-forward-probability}
\end{equation}
where $E(c)=|\{l:n_l=0\}|$ is the number of empty components before the split.
The reverse move is deterministic merge, so
\[
    \log q(c\mid c')=0.
\]
The split acceptance probability is
\begin{equation}
    \alpha_{\mathrm{split}}
    =
    \min\left\{1,
    \exp\left[
    \log\pi(c'\mid z)-\log\pi(c\mid z)
    -\log q(c'\mid c)
    \right]
    \right\}.
    \label{eq:split-acceptance}
\end{equation}

\subsection{Merge proposal}

Suppose $c_u=a$, $c_v=b$, and $a\neq b$.  Let
\[
    A=\{q:c_q=a\},
    \qquad
    B=\{q:c_q=b\}.
\]
The proposal merges $B$ into $A$:
\[
    c_q'=a,
    \qquad q\in A\cup B.
\]
Thus
\[
    \log q(c'\mid c)=0.
\]
The reverse proposal is a split of $A\cup B$ back into $A$ and $B$ using the
same anchors $u$ and $v$.  Let $p_q^{\mathrm{rev}}$ be the normalized
probability of assigning $q$ to its original side during this reverse
restricted Gibbs scan.  Then
\begin{equation}
    \log q(c\mid c')
    =
    -\log E(c')
    +
    \sum_{q\in(A\cup B)\setminus\{u,v\}}\log p_q^{\mathrm{rev}},
    \label{eq:merge-reverse-probability}
\end{equation}
where $E(c')$ is the number of empty components after the merge.  The merge
acceptance probability is
\begin{equation}
    \alpha_{\mathrm{merge}}
    =
    \min\left\{1,
    \exp\left[
    \log\pi(c'\mid z)-\log\pi(c\mid z)
    +\log q(c\mid c')
    \right]
    \right\}.
    \label{eq:merge-acceptance}
\end{equation}

\section{Posterior cluster summaries}

Because labels are not identifiable, cluster summaries should be
label-switching invariant.  The primary summary is the posterior similarity
matrix
\begin{equation}
    C_{qr}
    =
    P(c_q=c_r\mid Y),
    \qquad q,r=1,\dots,P.
\end{equation}
During MCMC, update online by
\begin{equation}
    C_{qr}^{\mathrm{raw}}
    \leftarrow
    C_{qr}^{\mathrm{raw}}+\ind(c_q=c_r)
\end{equation}
after burn-in.  At the end, divide by the number of saved iterations.

Also monitor the trace of the occupied-component count
\begin{equation}
    K_+^{(t)}=|\{l:n_l^{(t)}>0\}|.
\end{equation}


\section{One full MCMC sweep}

\begin{algorithm}[H]
\caption{One MCMC sweep: LSIRM update and collapsed label update}
\small
\begin{algorithmic}[1]
\State \textbf{Input:} data $Y$, layer maps $q\mapsto(\ell(q),j(q))$, hyperparameters, proposal scales, $K^\star$, and $M_{\mathrm{SM}}$.
\State Update layer-specific nuisance parameters $\xi_\ell$ (e.g.\ $\gamma_\ell$, $\sigma_0^2$, $\lambda_{ij}^{(\ell)}$ for continuous layers; $\{\kappa_j\}$ for count layers via \eqref{eq:kappa-prior} log-normal MH per item), if present.
\For{$i=1,\dots,n$}
    \State Propose $a_i'$ and accept/reject using \eqref{eq:a-mh-ratio} (shared across layers; prior variance fixed at $1$).
\EndFor
\For{$\ell=1,\dots,L$}
    \For{$i=1,\dots,n$}
        \State Propose $\alpha_i^{(\ell)\prime}$ and accept/reject using \eqref{eq:alpha-mh-ratio}.
    \EndFor
    \For{$j=1,\dots,P_\ell$}
        \State Propose $\beta_j^{(\ell)\prime}$ via \eqref{eq:beta-mh-ratio} for non-ordinal layers, or each threshold $\beta_{j,k}^{(\ell)\prime}$ via \eqref{eq:beta-thr-mh-ratio} for ordinal layers (rejecting any proposal that violates the ordering).
        \State Let $q=q(\ell,j)$ and propose $b_j^{(\ell)\prime}$.
        \State Accept/reject $b_j^{(\ell)\prime}$ using one of: \eqref{eq:b-mh-ratio} (current-label conditional), \eqref{eq:b-mh-ratio-inflated} ($\alpha$-inflated), or \eqref{eq:b-mh-ratio-marginal} (marginalized over $c_q,\rho,\mu,\Sigma$).
    \EndFor
    \State Update $\sigma_{\alpha,\ell}^2$ using \eqref{eq:sigma-alpha-update} with layer-specific hyperparameters $(A_{\alpha,\ell},B_{\alpha,\ell})$.
\EndFor
\State Construct global item positions $z_q=b_{j(q)}^{(\ell(q))}$ for $q=1,\dots,P$.
\For{$q=1,\dots,P$}
    \State Remove item $q$ from its current cluster.
    \State Compute NIW leave-one-out predictive weights in \eqref{eq:c-update} for $l=1,\dots,K^\star$.
    \State Sample $c_q$ from the normalized weights.
\EndFor
\end{algorithmic}
\end{algorithm}

\begin{algorithm}[H]
\caption{One MCMC sweep: split--merge, mixture-parameter update, and storage}
\small
\begin{algorithmic}[1]
\For{$m=1,\dots,M_{\mathrm{SM}}$}
    \State Select two anchors $u,v$ uniformly at random.
    \If{$c_u=c_v$}
        \State Attempt a restricted-Gibbs split move and accept/reject using \eqref{eq:split-acceptance}.
    \Else
        \State Attempt a merge move and accept/reject using \eqref{eq:merge-acceptance}.
    \EndIf
\EndFor
\For{$l=1,\dots,K^\star$}
    \State Compute NIW posterior parameters $(\kappa_l,\nu_l,m_l,S_l)$ from \eqref{eq:niw-posterior-params}.
    \State Draw $\Sigma_l\mid z,c$ using \eqref{eq:sigma-posterior}.
    \State Draw $\mu_l\mid \Sigma_l,z,c$ using \eqref{eq:mu-posterior}.
\EndFor
\State Draw $S_0\mid\{\Sigma_l\}$ using \eqref{eq:S0-gibbs} and refresh $\log|S_0|$ cache.
\State Optionally draw $\rho\mid c$ using \eqref{eq:rho-posterior}.
\If{iteration is after burn-in}
    \State Store $c$, $K_+$, $z$, and selected LSIRM parameters.
    \State Update posterior similarity counts $C_{qr}^{\mathrm{raw}}\leftarrow C_{qr}^{\mathrm{raw}}+\ind(c_q=c_r)$.
\EndIf
\end{algorithmic}
\end{algorithm}

\end{document}
